\title{Large Language Models Can Automatically Engineer Features for Few-Shot Tabular Learning}

\begin{abstract}
Large Language Models (LLMs) possess impressive capabilities for solving complex reasoning challenges and novel tasks, making them highly promising for tabular learning, which is fundamental to a variety of practical applications. We introduce an innovative in-context learning approach, \textsf{FeatLLM}, in which LLMs serve as feature engineers to generate a dataset formatted for optimal tabular prediction. These derived features inform class probability estimates through a straightforward downstream machine learning method—such as linear regression—resulting in effective few-shot learning. Our proposed \textsf{FeatLLM} system exclusively employs this basic predictive model and the constructed features at inference, distinguishing itself from other LLM-based methods by removing the requirement to query the LLM for every individual sample during inference. Additionally, \textsf{FeatLLM} only needs API access to LLMs and addresses constraints related to prompt length. Experiments conducted on diverse tabular datasets reveal that \textsf{FeatLLM} formulates high-quality predictive rules and significantly (by 10\% on average) outperforms competing methods including TabLLM and STUNT.
\end{abstract}

\section{Introduction}
Recent advances in Large Language Models (LLMs) have revealed their exceptional proficiency in generating context-appropriate outputs for tasks not seen during training, all without specialized fine-tuning~\cite{creswell2022selection,imani2023mathprompter,taylor2022galactica}. Leveraging vast textual data, the inherent prior knowledge embedded in LLMs can frequently replace the necessity for expert knowledge across various fields~\cite{Genomes2021,singhal2023large,wilcox2003role}, supporting automation in numerous scenarios such as dialogue processing~\cite{finch2023leveraging} and code repair~\cite{fan2023automated}. Notably, by including task details and a handful of examples within prompts, LLMs are able to contextualize assignments and highlight pertinent features and instances~\cite{brown2020language,zhao2023survey}, showcasing their potential for sophisticated data analysis and their suitability as automated feature engineers.

Efforts to leverage the reasoning ability and accumulated knowledge of LLMs have recently extended into the realm of tabular learning. For example, background information—for instance, correlations between age and certain health conditions—can enhance tasks like disease prediction. In such settings, recent literature defines both task and feature context in natural language, transforms tabular data into serialized forms, and incorporates this material into LLMs~\cite{dinh2022lift,wang2023anypredict}. Approaches generally proceed with in-context learning~\cite{nam2023semi} or efficient parameter fine-tuning~\cite{dinh2022lift,hegselmann2023tabllm}. Prior knowledge encoded in LLMs has consistently demonstrated benefits, exceeding traditional tabular learning benchmarks especially when training data is scarce~\cite{hegselmann2023tabllm}.

Current techniques for tabular learning utilizing LLMs exhibit several shortcomings. Achieving end-to-end predictions typically mandates at least one LLM inference per instance, resulting in considerable computational overhead. Furthermore, the majority of methods depend on LLM fine-tuning to reach competitive accuracy, which restricts their practical deployment, especially when LLMs lack open training infrastructure or tuning APIs—a significant issue as recent high-performing LLMs are increasingly accessible only through inference APIs~\cite{anil2023palm,openai2023gpt4}. Most existing solutions also fail to handle extended prompt lengths effectively~\cite{liu2023lost}. As the number of columns in tabular data increases and the corresponding serialized text expands, these techniques become impractical or suffer performance degradation. Collectively, these issues pose major obstacles to applying LLM-based tabular models in practical settings.

We propose an alternative approach that moves away from utilizing LLMs purely for end-to-end prediction. Instead, our method repurposes LLMs for feature construction, harnessing their extensive background knowledge. Concretely, we put forward a new in-context learning methodology named \textsf{FeatLLM}, which endeavors to reveal the underlying “criteria” guiding LLM predictions. For example, in disease diagnosis tasks, an LLM may articulate explicit rules that link feature patterns to disease identification. These extracted criteria are then used to synthesize new features, replacing traditional ones for task resolution. By delegating feature generation to the LLM, we can employ a streamlined downstream model—once meaningful features have been established, inference does not require complex models, resulting in improved inference speed. By exploiting the in-context learning capability of LLMs, feature extraction can be accomplished solely through incorporating demonstrations within prompts, obviating the need for any LLM retraining. Additionally, ensemble-inspired strategies such as feature bagging~\cite{breiman2001random} are integrated to alleviate problems associated with large prompt sizes.

\textsf{FeatLLM} decomposes the process of engineering novel features within the input prompt into two distinct stages: (1) analyzing the problem to deduce how features correlate with target outcomes, and (2) utilizing the previously inferred knowledge along with limited-shot examples to generate discriminative rules for class separation. By imposing this step-wise logical framework, LLMs are guided to disregard irrelevant features during the rule formulation phase. The resulting rules are converted to executable code and applied to the dataset, yielding binary indicators for each sample according to rule compliance. Subsequently, a simple model is constructed to infer class probabilities using these transformed features. Robustness is enhanced through an ensemble technique, involving repeated feature generation and feature bagging; varying the number of selected features and samples in each iteration helps to avoid prompt length limitations. Notably, \textsf{FeatLLM} operates without the need for model retraining and incurs minimal inference overhead, thus enabling practical adoption of LLMs in a variety of real-world contexts.

In our empirical study, \textsf{FeatLLM} is assessed on 13 different tabular datasets under low-shot learning scenarios, demonstrating consistently strong and stable performance. The results indicate that LLMs effectively harness both pre-existing domain knowledge and information directly gleaned from the data to construct valuable features. Our method surpasses several state-of-the-art few-shot learning baselines across multiple experimental configurations. The implementation is accessible through an anonymized GitHub repository at~\texttt{https://github.com/Sungwon-Han/FeatLLM}.

\section{Related Work}

\subsection{Few-Shot Learning with Tabular Data} The development of few-shot learning methods has made it possible to derive broadly applicable feature representations even when only limited label information is available~\cite{chen2018closer}. While few-shot learning approaches were initially tailored for visual data, they have been adapted successfully to various types of data, such as text, audio, and more recently, tabular formats~\cite{majumder2022few,nam2023stunt,schick2021exploiting}. When working with small labeled datasets, there is an inherent risk of models capturing non-generalizable patterns, i.e., spurious correlations~\cite{bartlett2021deep}. In response, contemporary strategies integrate supplementary information or embed prior domain-specific knowledge to guide model learning towards meaningful inductive biases. Strategies include leveraging abundant unlabeled tabular datasets and applying appropriate data augmentation techniques within a semi-supervised framework~\cite{ucar2021subtab,yoon2020vime}, as well as initializing models via unsupervised pre-training that does not rely on task-specific labels~\cite{somepalli2021saint}. Such pre-training may utilize self-supervised objectives like masking or corrupting segments of the input and reconstructing them~\cite{arik2021tabnet,majmundar2022met}, or using augmented data to create contrastive learning paradigms~\cite{bahri2022scarf,chen2020simple}. Nevertheless, the diverse and complex structure of tabular data presents difficulties in crafting augmentation schemes that are universally beneficial, leading to limited effectiveness for these approaches in few-shot settings~\cite{nam2023stunt}.

Recent advances have involved pre-training models on an extensive array of tabular datasets that may not be directly linked to the final task, followed by transfer learning techniques aimed at the specific task of interest~\cite{zhang2023generative,levin2022transfer,zhu2023xtab}. For example, TransTab~\cite{wang2022transtab} introduces transformer architectures that incorporate semantic information from column names, thereby enriching their learning capacity beyond just table cell values. The accumulated knowledge from these varied table structures and column semantics enables effective zero-shot and few-shot predictions for new tabular tasks. Conversely, TabPFN~\cite{hollmann2022tabpfn} generates synthetic datasets from mixed distribution types to pre-train a Bayesian neural network, enabling it to perform inference directly with the target task's samples—both training and test—without additional fine-tuning. This accumulation of prior knowledge from multiple datasets has led TabPFN to exceed the performance of classic tree-based methods and demonstrated substantial gains in few-shot contexts.

\subsection{Language-Interfaced Tabular Learning} A further promising direction for exploiting prior knowledge involves the use of large language models (LLMs)~\cite{anil2023palm,openai2023gpt4}. Through training on vast, heterogeneous text collections, LLMs have attained the ability to generalize to previously unseen tasks and domains~\cite{brown2020language}. Recent research initiatives have begun to utilize the latent knowledge embedded within LLMs for tabular learning by serializing tabular data as text and introducing them into LLM prompt inputs~\cite{dinh2022lift,hegselmann2023tabllm,wang2023anypredict}. By incorporating tabular content, feature and task descriptions within prompts, LLMs can discern functional relationships among features and tasks to execute inference. Progress in this area includes specialized model training via parameter-efficient adaptation methods such as LoRA~\cite{hu2021lora} or IA3~\cite{liu2022few}, alongside in-context learning approaches, which involve embedding few-shot example instances within the prompts, bypassing further training~\cite{dinh2022lift,hegselmann2023tabllm,nam2023semi,slack2023tablet}.

While the methods discussed above have demonstrated success in few-shot tabular learning, their reliance on routing inference exclusively through the LLM limits practicality for industrial deployment. To address this, our work proposes a paradigm shift from the traditional approach where every sample’s prediction is handled end-to-end by an LLM acting as a black box. Instead, our approach leverages the LLM for directly deducing class-specific rules or conditions. \textsf{FeatLLM} translates these learned rules into program code, enabling the synthesis of novel features. This allows for prediction without engaging the LLM at inference time, using in-context learning to induce rules by drawing on prior knowledge and the information found in a handful of examples.

\section{Methods}
\textsf{FeatLLM} focuses on identifying and extracting the underlying ``rules"—criteria corresponding to each label class—using the available limited examples, as illustrated in Figure~\ref{fig:main_model}. These rules are extracted through the LLM, and subsequently transformed into binary programmatic features that indicate whether each rule holds true for a given data point. The generated binary feature vector is then combined with a set of non-negative trainable weights via a dot product, producing a class-wise likelihood score. The learning process adjusts these weights, thus quantifying the influence of each feature on class membership (details provided in Section~\ref{Sec:training}). This methodology employs multiple iterations with bagging to generate diverse ensembles, which are later aggregated for final predictions. The upcoming sections elaborate on the problem definition and provide a comprehensive overview of each methodological component.

\vspace*{-0.12in}
\paragraph{Problem formulation.} Consider a tabular dataset composed of $N$ annotated examples, represented as $\mathcal{D} = \{(\mathbf{x}^i, \mathbf{y}^i)\}_{i=1}^{N}$. Each instance $\mathbf{x}^i$ consists of $d$ attributes, making it a vector in $d$ dimensions. These features correspond to descriptive names in natural language, denoted as $F = \{f_j\}_{j=1}^{d}$, for which separate definitions are available. In the context of classification, each label $\mathbf{y}^i$ is assigned from a predefined category set. For the $k$-shot learning scenario, a subset of $k$ samples (where $k < N$) is drawn at random from $\mathcal{D}$ to train the algorithm.

\subsection{Prompt Design for Extracting Rules}
\label{Sec:prompt_design}
To improve the fidelity of rule extraction by LLMs, we steer the reasoning process to resemble the methodology an expert would apply to tabular prediction challenges. The prompt strategy is structured around three main sections, detailed in Fig.~\ref{fig:prompt_for_rule} and Appendix~\ref{sec:example_rule}:

\vspace*{-0.12in}
\paragraph{Basic information description.} This section details all critical data necessary for tackling the task (see orange annotation in Fig.~\ref{fig:prompt_for_rule}), including the question framing with explicit label guidance, definitions and descriptors for all input features, and illustrative demonstrations. The question is phrased as a query related to the task (e.g., ``Does this patient's myocardial infarction complications data indicate chronic heart failure? Yes or no?"). Additional task description examples are available in Appendix~\ref{sec:task_desc}. Feature descriptions specify the data type and may optionally provide formal definitions; for categorical attributes, possible values are enumerated. As part of the demonstrations, select training examples are converted into textual format along with corresponding true labels. The serialization follows established conventions from prior works~\cite{dinh2022lift,hegselmann2023tabllm}:

\begin{align}
&\text{Serialize}(\mathbf{x}^i,\mathbf{y}^i, F) = \nonumber \\ 
&``\ f_1\ \text{is}\ \mathbf{x}^i_1.\ ... \ f_d\  \text{is}\  \mathbf{x}^i_d.\ \ \text{Answer:}\  \mathbf{y}^i\ ”
\end{align}

\vspace*{-0.12in}
\paragraph{Reasoning instruction.} Our goal is to improve the reasoning capabilities of the LLM by incorporating explicit reasoning instructions (highlighted in blue in Fig.~\ref{fig:prompt_for_rule}). The reasoning guidance starts with a sentence modeled after the chain-of-thought methodology~\cite{wei2022chain}. Subsequently, we segment the inference procedure into two distinct phases. During the initial phase, the LLM is prompted to discern the causal links or inclinations connecting features to the task description, independent of any demonstration examples. Instead, this stage leverages the model's general understanding or common sense to allow thorough utilization of its prior knowledge regarding the problem. In the following phase, the LLM integrates the example demonstrations and insights derived from the first stage to construct rules corresponding to each class. This bifurcated reasoning strategy serves to minimize the identification of superficial correlations in extraneous columns, thereby encouraging the model to concentrate on salient features. To avoid both the risk of underfitting and the complications from an unwieldy prompt due to excessive rule extraction, we fix the number of rules extracted per class at 10\footnote{Section~\ref{sec:analysis} offers further discussion on the selection of rule quantities.}. Moreover, when generating rules, the LLM is instructed to reference the data type of each feature indicated in the basic information description, ensuring compatibility of rule logic with feature types.

\vspace*{-0.12in}
\paragraph{Response instruction.} To enable efficient parsing and utilization of the rules generated, we provide explicit instructions on response formatting to the LLM (yellow text in Fig.~\ref{fig:prompt_for_rule}). These prompts detail the format requirements for responses in relation to each answer class (i.e., response format) as well as the prescribed structure for individual rules (i.e., rule format). Such instructions empower the LLM to adjust the formatting to suit the feature type appropriately. In the absence of further directives, the LLM can interlink multiple rules employing logical operators such as AND or OR. A sample output is documented in Appendix~\ref{sec:outcome-rule}.

\subsection{Inferring Class Likelihood via Rules}
\label{Sec:training}

\paragraph{Rule parsing for automated feature extraction.} In this stage, we transform the previously obtained rules into binary features for each target class. Specifically, for every class, these features signal whether a given sample meets the conditions of that class’s rule set, using a binary indicator ('0' or '1'). Such features serve as inputs for determining the class membership likelihood of samples. Due to the LLM’s reliance on natural language for rule construction, a dedicated parsing mechanism is essential to reliably convert rules into features. Despite imposing constraints on the output structure and rule formatting to streamline this process, the LLM may still introduce syntactic inconsistencies or deviations~\cite{kovavc2023large}, such as replacing mathematical symbols like ``$>$” or ``$<$” with textual equivalents like ``greater than" or ``less than," thereby complicating parsing.

To mitigate difficulties arising from the unpredictable syntax, we adopt an approach that employs the LLM itself rather than constructing elaborate programmatic parsers. The LLM’s proficiency in comprehending meaning despite syntactic alterations, alongside its capability to produce concise code in response to task instructions (as it is trained on code-based datasets), allows us to exploit these strengths. By providing prompts that contain the function name, details of input and output, as well as the identified rules, we delegate the feature-parsing task to the LLM. Illustrative prompts and their corresponding results can be found in Figures~\ref{fig:prompt_for_parsing} and~\ref{sec:example_parsing}-\ref{sec:example_parse_outcome} in the Appendix. The Python \texttt{exec()} function is used to execute the code generated by the LLM, completing the data transformation.

\vspace*{-0.12in}
\paragraph{Estimating likelihood for class assignment.} With the collection of binary features derived from class-specific rules, a straightforward approach for estimating the likelihood that a sample belongs to a certain class is to aggregate the satisfied rule indicators for that class, effectively summing the binary features per class. Nonetheless, individual rules and features may differ in predictive value, suggesting the need to learn their relative weights from training data. We propose leveraging standard machine learning (ML) models to discern these feature importances within each class’s binary indicators. For instance, a bias-free linear model may be utilized for this purpose. The binary features generated from sample $\mathbf{x}^i$ for class $k$ are represented as $\mathbf{z}_k^i \in \{0, 1\}^R$, where $R$ denotes the count of rules per class and $c$ the total number of classes. The class probabilities $\mathbf{p}^i$ are then calculated as follows:

\begin{align}
\text{logit}_k^i &= \max(\mathbf{w}_k, 0) \cdot \mathbf{z}_k^i, \nonumber \\
\mathbf{p}^i &= \text{Softmax}({[\text{logit}_{1}^i, ..., \text{logit}_{c}^i]}). \label{eq:infer_prob}
\end{align}

In this context, $\mathbf{w}_k$ denote the parameters of the linear model that are subject to optimization and reflect the relevance of individual features. By constraining these weights to lie within a positive domain, we prevent the features derived from the LLM from negatively influencing the class probability during the learning phase. The logit outputs are mapped to class probabilities by applying the softmax transformation. Training involves minimizing the cross-entropy loss with respect to the weights $\mathbf{w}_k$, based on a limited set of annotated samples.

\vspace*{-0.12in}
\paragraph{Ensembling with bagging.} To obtain the final prediction, we repeat the entire procedure multiple times, thereby generating an ensemble of models. Let the trained weights from $T$ repetitions be denoted as $\{\mathbf{w}[t]\}_{t=1}^T$. The final class probability for $\mathbf{p}^i$ is calculated as the mean of the probabilities yielded by each instance $\mathbf{w}[t]$, following Eq.~\ref{eq:infer_prob}.

Randomness is incorporated in the rule generation for the ensemble by assigning a sufficiently elevated temperature parameter during LLM inference. Additionally, the sequence of few-shot examples is shuffled since prior findings indicate that the order may influence LLM output~\cite{lu2022fantastically}\footnote{Further evaluation of rule diversity is presented in Appendix~\ref{sec:diversity}}. Bagging is applied to sample subsets of features or data points for each iteration, thus harnessing prediction diversity for improved robustness—a strategy especially advantageous when the dataset involves a large number of features or training examples. This ensemble strategy confers two primary benefits. Firstly, erroneous rules produced by the LLM in some runs may be offset by accurate rules from other runs, thereby increasing the system’s resilience. This is related to the self-consistency property of LLMs~\cite{wang2022self}, as rules that are frequently generated across several repetitions are often more reliable. Secondly, the ensemble mechanism mitigates restrictions imposed by LLM prompt length. For tabular datasets exceeding the prompt capacity, bagging enables prompt reduction, maintaining reliable model performance. While a single rule set may not adequately describe relationships among all features, combining multiple weak learners across different trials via ensembling diminishes the effects of any missing feature or instance coverage in individual runs.

\section{Experiments}
We assess the performance of \textsf{FeatLLM} on a range of tabular datasets and systematically analyze its mechanisms. Owing to limitations in space, supplemental experiments—such as evaluations with different LLM variants, in-depth qualitative analysis, and further studies—are deferred to the Appendix.

\subsection{Performance Evaluation}
\paragraph{Datasets.}
Our empirical studies involve 13 datasets designed for either binary or multi-class classification tasks. These are: (1) Adult~\cite{asuncion2007uci}, which predicts individuals earning over \$50,000 per year; (2) Bank~\cite{moro2014data}, which determines whether a client subscribes to a term deposit; (3) Blood~\cite{yeh2009knowledge}, aimed at forecasting if donors return for another donation; (4) Car~\cite{kadra2021well}, which assesses car quality; (5) Communities~\cite{misc_communities_and_crime_183}, which estimates crime rates in specific areas; (6) Credit-g~\cite{kadra2021well}, for the prediction of good versus bad credit risks; (7) Diabetes\footnote{\texttt{kaggle.com/datasets/uciml/pima-indians-diabetes-database}}, for diagnosing diabetes in individuals; (8) Heart\footnote{\texttt{kaggle.com/datasets/fedesoriano/heart-failure-prediction}}, to predict the likelihood of coronary artery disease; and (9) Myocardial~\cite{misc_myocardial_infarction_complications_579}, used to determine the presence of chronic heart failure.

Additionally, the study encompasses novel or synthetically generated datasets, which have rarely been discussed and are absent from LLM pre-training corpora: (10) Cultivars, which forecasts soybean cultivar grain yield\footnote{\texttt{archive.ics.uci.edu/dataset/913}}; (11) NHANES, which classifies individuals into age categories based on their records\footnote{\texttt{archive.ics.uci.edu/dataset/887}}; (12) Sequence-type, which involves classifying a numerical sequence as either arithmetic, geometric, Fibonacci, or Collatz; and (13) Solution-mix, which predicts whether a compound’s concentration derived from four constituent solutions exceeds 0.5. The first two of these have been contributed to the UCI Machine Learning Repository post-September 2023, ensuring that they do not appear in the LLM API (gpt-3.5-turbo-0613) pre-training data. The last two are generated synthetically to prevent their content from being memorized by the LLM used in this evaluation.

A diversity of dataset scales and complexities is represented, further detailed in Table~\ref{Tab:basic-desc}.
Each tabular dataset is accompanied by explicit attribute names and descriptions. Some datasets, such as `Communities' and `Myocardial', contain over a hundred attributes, which presents additional difficulties given LLM context length constraints.

\vspace*{-0.12in}
\paragraph{Baselines.}
For comparative analysis, we include nine baseline methods. The first set features three traditional algorithms for tabular classification: (1) Logistic Regression (LogReg), (2) XGBoost~\cite{chen2016xgboost}, and (3) RandomForest~\cite{ho1995random}. Following are two methods leveraging unlabeled datasets: (4) SCARF~\cite{bahri2022scarf} and (5) STUNT~\cite{nam2023stunt}, though acquiring ample unlabeled data can be problematic in practice. Another competing approach is (6) TabPFN~\cite{hollmann2022tabpfn}, which utilizes synthetically generated, diversely distributed data for pretraining. The final three utilize LLMs: (7) In-context learning (In-context)~\cite{wei2022emergent}, (8) TABLET~\cite{slack2023tablet}, and (9) TabLLM~\cite{hegselmann2023tabllm}. In-context learning incorporates a small selection of labeled examples into the LLM prompt without parameter updates. TABLET augments the prompt with external hints from classifiers based on rule-sets and prototypes, thereby enhancing inference reliability. TabLLM employs parameter-efficient techniques such as IA3~\cite{liu2022few}, fine-tuning the LLM with a few labeled examples.

\vspace*{-0.12in}
\paragraph{Implementation details.}
\textsf{FeatLLM} adopts GPT-3.5 as its main LLM architecture, although its framework supports flexibility in LLM selection (refer to Appendix~\ref{sec:backbones} for experimental outcomes with various architectures). The LLM inference process utilizes a temperature parameter of 0.5 and maintains the API's standard top-p setting of 1. The ensemble size is established at 20, and 10 rules are used for extraction. The influence of these hyper-parameters is examined in Figure~\ref{fig:hyperparameter2} and discussed further in Appendix~\ref{sec:hyper-parameter}. 

For the linear model, we optimize using Adam with a learning rate of 0.01, running the training procedure for 200 epochs. $k$-fold cross-validation is employed to choose the best epoch. When duplicating baseline methods, GPT-3.5 serves for in-context learning, while T0~\cite{sanh2022multitask} is used as the LLM backbone in TabLLM, in accordance with prior configurations. It should be highlighted that TabLLM limits LLM choices to those with open checkpoints, thus GPT-3.5 is excluded from these benchmarks. 

Parameter tuning for standard machine learning algorithms (LogReg, XGBoost, and RandomForest) is performed through Grid-search paired with $k$-fold cross validation. The value of $k$ is set at either 2 or 4, ensuring each training set has at least one sample per class. Additional implementation specifics are found in Appendix~\ref{sec:baselines}.

\vspace*{-0.12in}
\paragraph{Results.}
Tables~\ref{tab:main1} and \ref{tab:main2} display comparative results for various datasets. Each experiment was conducted three times, each with a distinct randomized split of training and test sets; the reported values represent the average AUC and its standard deviation across these runs. Across baseline models, \textsf{FeatLLM} typically achieves the highest or second-highest AUC scores. Our approach, using only 4-shot learning, matches the performance of traditional tabular baselines that operate with 32 shots, as illustrated in Fig.~\ref{fig:compares}. Although STUNT and TabPFN obtained notable improvements on some datasets, their overall advantage over standard machine learning approaches was modest. Additionally, LLM-based methods like in-context learning, TABLET, or TabLLM could not handle inference tasks for datasets with large feature sets. In contrast, our model, which integrates bagging and ensemble methodologies, reliably delivers strong results even for datasets containing many features. It is also worth mentioning that applying the bagging and ensemble strategy to existing LLM-oriented techniques is theoretically possible, but it would require twice the computational effort per inference sample, resulting in significant computational overhead that renders such modifications impractical.

\subsection{Analysis \& Discussion}
\label{sec:analysis}
\paragraph{Ablation study.} We conduct ablation experiments to evaluate the significance of key components in our method, considering the following: (1) \textsf{-Tuning}: excludes weight tuning in the linear model, instead directly aggregating binary features to estimate logits; (2) \textsf{-Ensemble}: removes ensemble operations; (3) \textsf{-Description}: omits feature description; (4) \textsf{-Reasoning}: skips Step-1 of reasoning instruction during rule generation.

Table~\ref{tab:ablation} presents the averaged impact of these ablations on AUC across all datasets. In every scenario, altering any component led to diminished performance. Notably, the value of weight tuning increases with higher shot counts, implying that ample data enables better identification of rule relevance, thereby amplifying the effect of tuning. Conversely, feature description and reasoning instruction have more pronounced effects when shot numbers are low, underlining the importance of leveraging LLMs' prior knowledge for performance gains. The ensemble strategy delivers persistent performance improvement irrespective of settings.

\vspace*{-0.12in}
\paragraph{Training \& inference time comparison.} We evaluate and contrast the durations required for training and inference across various approaches. These tests utilize the Adult dataset, with a training set comprising 16 instances. For most methods, a single A100 GPU is deployed, while TabLLM employs four A100 GPUs to facilitate model parallelism. The runtime statistics are presented in Table~\ref{tab:cost}. Significantly, \textsf{FeatLLM} achieves inference times that are nearly equivalent to those of standard baseline techniques. In comparison, methods based on LLMs—including In-context learning, TABLET, and TabLLM—require much greater inference time. With respect to training time, \textsf{FeatLLM} is on par with other LLM-driven models; its requirement of only 30 API calls constitutes a manageable computational demand, and these queries are amenable to parallel execution.

\vspace*{-0.12in}
\paragraph{Quality of features generated by \textsf{FeatLLM}.}
To assess the usefulness of the rule-based features produced by \textsf{FeatLLM} for real-world tasks, we undertake a straightforward study. We evaluate the mean AUROC score between the newly created features and the target labels, and calculate the proportion of features whose AUROC exceeds 0.5 across each dataset. Table~\ref{tab:quality} presents these findings. The evidence demonstrates that most of the derived rules are pertinent to the target variable, thus offering valuable support in tackling the task at hand (see the column ``Before tuning"). Additionally, when tuning a linear model using available data, features with marginal relevance (namely, those with associated weights below a specified threshold) are systematically excluded (refer to the column ``After tuning"). The threshold, set at 0.1, is determined based on the distribution of learned weights.

\vspace*{-0.12in}
\paragraph{Impact of incorporating spurious correlations.} We investigated the ability of several methods to eliminate noisy spurious correlations under a low-shot scenario. For this purpose, we utilized the Heart dataset and randomly introduced columns from the Adult dataset that bear no relevance to the Heart dataset's prediction task. We assessed model performance as the quantity of such irrelevant columns gradually increased. To ensure a fair evaluation, all approaches were provided with 8 labeled samples and did not leverage any unlabeled data, thereby excluding algorithms such as SCARF and STUNT from this comparison.

Figure~\ref{fig:spurious} illustrates how spurious correlations affect model accuracy. Among all methods, \textsf{FeatLLM} was the least affected in terms of performance deterioration caused by irrelevant features. This resilience may stem from the framework's explicit use of prior knowledge to align features with the task during rule extraction, thus discouraging the formation of rules based on nonessential columns and resulting in greater robustness. Notably, rules formed from noisy columns constituted only about 1-2\% of the total rule set. Nonetheless, we found that \textsf{FeatLLM} may still capture spurious associations and misinterpret causal dependencies when columns from highly similar data sources are arbitrarily appended (refer to Appendix~\ref{sec:spurious-appendix} for details).

\vspace*{-0.12in}
\paragraph{Hyper-parameter analysis}
\textsf{FeatLLM} is governed by two key hyper-parameters: the ensemble count and the number of rules per class extracted during each LLM invocation. To understand their influence on overall performance, we performed a series of experiments. Our findings indicate that increasing the ensemble count steadily enhances performance, but this improvement plateaus near 20 ensembles (refer to Fig.~\ref{fig:appendix-hyperparameter1} in the Appendix). Regarding the number of rules, although differences were not statistically substantial, selecting 10 rules appeared optimal when balancing prompt size against predictive accuracy (see Fig.~\ref{fig:hyperparameter2}). We hypothesize that adding more rules raises input dimensionality, which contributes additional information but also risks overfitting in scenarios with few examples.

\section{Conclusion}
We introduce \textsf{FeatLLM}, a method that raises the bar for LLM-driven few-shot tabular learning. Rather than merely leveraging LLMs for sample-level prediction, \textsf{FeatLLM} adopts an approach where the LLM generates predictive criteria, functioning as a feature engineering agent. Through rule generation and an ensemble scheme via bagging, \textsf{FeatLLM} minimizes inference requirements, operates exclusively with API access—requiring no model training—and sidesteps limitations related to feature dimensionality. Consequently, \textsf{FeatLLM} demonstrates strong predictive capabilities and emerges as an advantageous strategy for efficient deployment of LLMs on practical tabular datasets.

Currently, \textsf{FeatLLM} is tailored for situations with few available samples. For future work, we aim to extend its scope to facilitate feature creation for larger datasets, experiment with alternative feature types beyond rules, and further enhance interpretability, thereby broadening its usability to diverse task domains.

\section{Broader Impact}
\textsf{FeatLLM} is developed to harness the accumulated knowledge within LLMs for tabular learning, aiming to surpass the limitations of conventional supervised techniques in data-scarce settings. By leveraging LLM prior knowledge, our work addresses the tendency towards overfitting when training data is limited, providing supplementary information to bolster learning efficacy. The framework is particularly suited for real-world applications in domains such as finance and healthcare, where annotation is often expensive or difficult, and pretrained LLMs may contain valuable domain-specific knowledge. As a critical future direction for broader application, it will be essential to investigate how societal biases and misinformation embedded in LLM prior knowledge affect predictive outcomes, since these pretrained insights can directly shape—and potentially outweigh—patterns learned from collected labeled data.

\end{document}